% Page 22 — Alban en sa « légende dorée »

\begin{center}
  {\Large\itshape Alban en sa « légende dorée »}
\end{center}
\vspace{1.5em}

\noindent
\lettrine{E}{ffectivement} la Madone est là. Le bon larron aussi. Les portes du paradis sont toutes ouvertes. La cour céleste est sur les dents. La Madone réclame son enfant. Le démon exige son client. Et la dispute entre notre « céleste avocate » et l'enfer déchaîné éclate avec fracas. Le démon rugit. La Vierge sourit. La discussion durerait encore si le diable, à bout d'arguments, n'avait proposé un pari. Alain de Nougaret appartiendra corps et âme à qui lancera le plus loin la plus grosse pierre. « D'accord – qu'on monte au pic du Montalet ».

\medskip

\noindent
\lettrine{A}{vec} ses ailes de chauve-souris, Satan a volé jusqu'au Montalet. Il arrache aux flancs de la colline une pierre énorme et de toutes ses forces, il lance dans l'espace ce bloc de rocher qui vient se lancer, droit comme un I, au chemin de la Bessière. Le Malin aurait-il gagné son pari ? Le pavé est en effet de taille respectable : colossal caillou qui est venu se placer à l'ouest du village de Notre Dame d'Alban. Quelle revanche pour Méphistophélès d'avoir lancé son pavé dans le jardin de la Vierge immaculée. Serait-ce une petite réplique vengeresse à Celui qui, le maudissant au paradis terrestre, lui assura que « la Femme lui écraserait la tête, tandis que lui, serpent venimeux, essaierait de la mordre au talon ? »

\medskip

\noindent
\lettrine{A}{lors} Satan se croit sûr de la victoire ; c'est que du Montalet ici, il y a plus qu'un port de fusil comme disent les paysans du coin. Et le pavé n'est pas une plume – C'est d'une chiquenaude que la Madone, elle, a lancé dans les rues son bloc de rocher, bien plus énorme que celui de Messire Satan. Et, comme un éclair, le formidable pavé va choir lui aussi aux portes de la cité confiée depuis des siècles et des siècles à Marie. Mais c'est 250 mètres plus loin que le caillou de Satan. La Vierge a donc gagné.

\medskip

\noindent
\lettrine{O}{uvrez-vous} donc, portes éternelles, et toutes grandes, pour Alain de Nougaret, lui aussi bon larron repentant. Et c'est entraîné par les mains maternelles de Notre Dame d'Alban et de partout, que le baron d'Alban entre dans les célestes parvis, cependant que les anges chantent à tue tête « Salve Mater Misericordiae », et que chérubins et séraphins redisent sans fin la louange de Celui dont c'est le propre d'avoir toujours pitié et de pardonner. Et c'est à l'aube des émancipations.
